% Chapter 8: OP3-C6C, P2b-B2/B3 and T1 exact original Fock operator algebra.
% The Hilbert--Schmidt derivation below is conventional and not separately Lean-qualified.
\chapter[Spin(9) on the Original Fock Algebra]{The Original Fock-Operator Action, Physical Adjoint, and Trace}
\label{ch:fock-action}

\section{The operator action is not yet a wavefunction action}
In Chapter~\ref{ch:clifford} we recovered a source-exact complex Clifford
algebra isomorphism
$\mathrm{Cl}_{48}^{\mathbb C}\cong\operatorname{End}_{\mathbb C}(\Fermion)$.
A spatial rotation acts on the original forty-eight real Majorana
labels through \eqref{eq:spin48-label-action}. A natural question is
whether we can transfer that transformation to the full original
Fock \emph{operator algebra}. That question differs from asking for a
continuous unitary operator $U_{\Fermion}(g)$ on Fock \emph{vectors}.
The distinction is not terminology: conjugation by $U$ forgets the
phase of $U$, and its effect on the Spin center must still be checked.

Let
\[
 \mathscr A:=\mathcal L_{\mathbb C}(\Fermion,\Fermion)
 \cong \operatorname{End}_{\mathbb C}(\Fermion),
 \qquad \Fermion=\mathbb C^{2^{24}}.
\]
The actual scalar field and adjoint in this definition are those of
\texttt{Fermion 2}. In particular, the star below is the adjoint with
respect to the \emph{original} finite-dimensional complex Hilbert
inner product, not an abstract Clifford reversal introduced for a
separate representation.

\section{Covariance of all forty-eight active Majorana generators}
Write $\Gamma:V_{48}\to_{\mathbb R}\mathscr A$ for the active original
real-linear normalized Majorana map. We distinguish this original map
from an earlier donor realization whose equivalence required a separate
source-label proof.

\begin{theorem}[Active original Majorana covariance; strict-qualified]
\label{thm:active-majorana}
There exists an exact source group homomorphism
\[
 \alpha:G_9\longrightarrow\operatorname{Aut}_{\mathbb C\text{-alg}}(\mathscr A)
\]
whose action agrees, for every $g\in G_9$ and original $v\in V_{48}$,
with the source forty-eight-label orthogonal representation:
\begin{equation}\label{eq:active-covariance}
 \alpha_g\bigl(\Gamma(v)\bigr)
   =\Gamma\bigl(R_{48}(g)v\bigr).
\end{equation}
The statement holds for an arbitrary real source vector, not merely
the forty-eight preferred coordinate axes.
\end{theorem}
\begin{proof}[Source-aligned proof argument]
The donor Clifford realization already carries the source Spin action,
with covariance against $R_{48}$. The new comparison identifies the
donor real-linear Majorana map with the \emph{active} original map.
It suffices to prove that equality on the source real-coordinate basis;
real-linearity then gives it on all of $V_{48}$. Substituting this equality
into the donor covariance theorem gives \eqref{eq:active-covariance}
without unfolding the enormous concrete finite Fock-operator sum.
The latter point is not merely an implementation convenience: it
certifies that the theorem refers to the original G1A-labelled CAR
operators, rather than a mathematically equivalent but unbound donor.
\end{proof}
\formalrecord{\path{OP3C6COriginalFockActiveGammaCovariance.lean},
\path{originalFockGammaRealLinear_eq_activeGammaRealLinear},
\path{source9Spin_originalFockOperatorAutomorphismHom_activeGamma};
strict-green original source and independent consumer
\href{https://github.com/etblink/math/actions/runs/42728268316}{\#42728268316}.}

\section{Preserving the genuine physical star}
A complex algebra automorphism need not preserve a chosen Hermitian
inner product. Thus \eqref{eq:active-covariance} by itself does not
show that rotations respect the \emph{physical adjoint}. The difference
was addressed through a separate sequence of small proof obligations:
from generator compatibility and scalar algebra-map behavior to the
result for arbitrary operators in the full finite Fock algebra.

\begin{theorem}[Original physical-star Spin action; strict-qualified]
\label{thm:fock-star-action}
The original source group acts by complex star-algebra automorphisms:
\begin{equation}\label{eq:fock-star-action}
 \alpha:G_9\longrightarrow
 \operatorname{Aut}_{\mathbb C,*}(\mathscr A),
 \qquad
 \alpha_g(A^*)=\alpha_g(A)^*,\quad
 \alpha_{gh}=\alpha_g\circ\alpha_h.
\end{equation}
Here $A^*$ is the original Hilbert adjoint on $\Fermion$.
\end{theorem}
\begin{proof}[Mathematical structure of the accepted source proof]
The original Clifford-to-Fock algebra equivalence of
Theorem~\ref{thm:clifford-equivalence} is surjective. Every original
Fock operator is therefore represented by an element of the source
complex Clifford algebra. The prerequisite star proof establishes
compatibility on the source Clifford generators, the scalar algebra
map, multiplication, and the linear operations needed for the
induction. Transporting those equalities through the original
algebra equivalence gives the adjoint-preserving statement for every
$A\in\mathscr A$. The accepted P2b-B3 stage packages the already
proved algebra action with that property as a genuine
\texttt{StarAlgEquiv}-valued group homomorphism. The group law is
inherited from the underlying source automorphism action, including
its order of composition.
\end{proof}
\formalrecord{\path{P2bB2OriginalFockGlobalPhysicalStar.lean},
\path{source9Spin_originalFockOperator_global_map_star};
\path{P2bB3OriginalFockStarAutomorphismAction.lean},
\path{source9Spin_originalFockOperatorStarAutomorphismHom}.
Strict-green P2b-B3 gate
\href{https://github.com/etblink/math/actions/runs/42764649026}{\#42764649026}.
The theorem is algebraic in the \emph{group parameter}; this gate
does not assert continuity $g\mapsto\alpha_g$ in a named topology.}

\section{A source trace, rather than a rescaled surrogate trace}
Finite dimensionality provides the ordinary complex trace
$\operatorname{Tr}_{\Fermion}$. It is important to use this particular
trace, rather than an unrelated functional obtained from a regular
Clifford representation: regular and irreducible traces have different
normalizations.

\begin{definition}[Original finite Fock trace]
Let $e:\operatorname{End}_{\mathbb C}(\Fermion)\to\mathscr A$
be the canonical finite-dimensional algebra equivalence between
linear maps and continuous linear maps. Define
\begin{equation}\label{eq:original-trace}
 \operatorname{Tr}_{\mathrm{orig}}(A)
  := \operatorname{Tr}_{\Fermion}(e^{-1}A),
  \qquad A\in\mathscr A.
\end{equation}
\end{definition}

\begin{theorem}[Trace invariance of the original action; strict-qualified]
\label{thm:original-fock-trace}
For every $g\in G_9$ and $A\in\mathscr A$,
\begin{equation}\label{eq:trace-invariant}
 \operatorname{Tr}_{\mathrm{orig}}(\alpha_g(A))
   =\operatorname{Tr}_{\mathrm{orig}}(A).
\end{equation}
\end{theorem}
\begin{proof}[Conjugation-free trace argument]
Transport $\alpha_g$ through $e$ to an algebra automorphism
$\beta_g=e^{-1}\alpha_g e$ of the finite-dimensional complex
endomorphism algebra. The standard trace-invariance theorem for algebra
equivalences of finite endomorphism algebras yields
$\operatorname{Tr}(\beta_g B)=\operatorname{Tr}(B)$ for all $B$.
Take $B=e^{-1}A$ and unfold \eqref{eq:original-trace}.
This proof does not require choosing a unitary conjugator or fixing
its phase on $\Fermion$.
\end{proof}
\formalrecord{\texttt{T1OriginalFockTraceInvariance.lean},
\texttt{originalFockTrace} and
\texttt{source9Spin\_originalFockTrace\_invariant}; strict-green
\href{https://github.com/etblink/math/actions/runs/42769361446}{\#42769361446}.
The source uses Mathlib's finite-dimensional
\texttt{LinearMap.trace\_map}.}

\section{The Hilbert--Schmidt consequence and what it does not give}
The three accepted statements---the group action, physical-star
preservation, and actual trace preservation---have a useful conventional
mathematical consequence that does \emph{not} require an arbitrary choice
of Fock-vector implementation.

\begin{proposition}[Algebraic Hilbert--Schmidt isometry; conventional deduction]
\label{prop:hs-isometry}
Equip $\mathscr A$ with the finite-dimensional Hilbert--Schmidt product
\begin{equation}\label{eq:hilbert-schmidt}
 \langle A,B\rangle_{\mathrm{HS}}
   :=\operatorname{Tr}_{\mathrm{orig}}(A^*B).
\end{equation}
Then the existing source automorphisms satisfy
\[
 \langle\alpha_g(A),\alpha_g(B)\rangle_{\mathrm{HS}}
   =\langle A,B\rangle_{\mathrm{HS}}
\]
for all $A,B\in\mathscr A$ and $g\in G_9$. Thus the \emph{already
constructed algebraic group homomorphism} acts by complex-linear
isometries on the Hilbert--Schmidt space.
\end{proposition}
\begin{proof}
Multiplicativity and star preservation give
\[
 \alpha_g(A)^*\alpha_g(B)
  =\alpha_g(A^*)\alpha_g(B)=\alpha_g(A^*B).
\]
Taking the original trace and using
Theorem~\ref{thm:original-fock-trace} proves the desired inner-product
identity. The group law is Theorem~\ref{thm:fock-star-action}.
\end{proof}
\editornote{This is a conventional algebraic corollary of accepted Lean
roots, not an independently qualified new Hilbert--Schmidt Lean declaration.
Strong/topological continuity of the source $G_9$ representation and
its physical significance are separate proof obligations.}

The dimension of this operator-valued fiber is
$\dim_{\mathbb C}\mathscr A=2^{48}$, while the dimension of the physical
Fock-vector fiber is $2^{24}$. Accordingly, even a genuine continuous
Hilbert--Schmidt action would not automatically be the original
wavefunction action. The input index of the operator fiber contributes
to representation-theoretic weights, a distinction that reappears in
the conditional low-oscillator exclusion estimates.

\section{The center and the missing implementer}
The original negative central element $z_9$ acts trivially on the
nine-dimensional \emph{vector} representation, but the separately
qualified source forty-eight-label homomorphism satisfies
\begin{equation}\label{eq:negative-center-48}
 R_{48}(z_9)=-I_{48}.
\end{equation}
Thus \eqref{eq:active-covariance} yields
$\alpha_{z_9}(\Gamma(v))=-\Gamma(v)$ on every original Majorana
label. Since the original Majoranas generate the full finite Fock
operator algebra, this is not the identity automorphism.

If an operator $U_{\Fermion}(z_9)$ implements this transformation by
conjugation on Fock vectors, its conjugation action agrees with that of
fermion parity. But a conjugation automorphism determines its implementer
only up to a complex scalar phase. One must establish the global
continuous group lift, the central sign/phase normalization, and the
required gauge behavior before asserting a particular vector action.
The tempting shortcut that sends the scalar $-1\in\mathrm{Cl}_9$
through a unital Clifford algebra homomorphism into scalar $-I$ on
Fock space would produce \emph{trivial} conjugation and therefore
cannot implement \eqref{eq:negative-center-48}. That rejected route
is important negative knowledge, not an obstruction to a correctly
constructed Spin group lift.

\begin{openproblem}[Original physical-vector Spin action]
Construct a continuous unitary representation
$U_{\Fermion}:G_9\to U(\Fermion)$ satisfying the exact original
Majorana covariance, correct source group law and central normalization.
Alternatively, prove that a suitably continuous Hilbert--Schmidt
operator-valued route suffices for an explicitly identified physical
spectral problem. In either case establish its precise relation to
the actual \texttt{FullL2 2} source representation and Hamiltonian.
\end{openproblem}

\section{Reader's checkpoint}
The original Fock \emph{operator} algebra is no longer a speculative
interface: it has a genuinely source-bound Spin action, physical-star
structure, and invariant trace. These are substantial foundations for
representation theory. They do not yet identify a continuous source
Spin action on original Fock vectors or its reducing sectors for the
self-adjoint physical Hamiltonian of Chapter~\ref{ch:hamiltonian}.
